\chapter{Kết Nối Thế Hệ (Bridging the Generations)}
\begin{center}
  \textit{\color{truyengold}``Xã hội lớn mạnh khi người già trồng cây mà họ biết mình không bao giờ được ngồi dưới bóng.''}\\[4pt]
  \textit{(A society grows great when old men plant trees whose shade they know they shall never sit in.)}\\[4pt]
  \small\color{truyenblue}--- Tục ngữ Hy Lạp ---
\end{center}
\ngancach

\section{Người Thầy Đặt Viên Gạch Cuối}
\begin{truyen}{Người Thầy Đặt Viên Gạch Cuối}{Câu chuyện về Ý Nghĩa Nghề Giáo}
\chuhoa{M}{ột} du khách đi ngang qua công trường xây dựng hỏi ba người thợ đang làm việc: ``Các anh đang làm gì vậy?''\\[4pt]
\textit{(A traveler passing a construction site asked three workers who were busy: `What are you doing?')}

Người thứ nhất cáu kỉnh: ``Tôi đang đặt gạch.'' Người thứ hai thở dài: ``Tôi đang xây tường.'' Người thứ ba ngẩng mặt lên, mắt sáng rỡ: ``Tôi đang xây một ngôi thánh đường.''\\[4pt]
\textit{(The first grumbled: `I am laying bricks.' The second sighed: `I am building a wall.' The third looked up, eyes shining: `I am building a cathedral.')}

Ba người cùng làm một việc --- nhưng chỉ một người thật sự hiểu ý nghĩa của công việc mình đang làm.\\[4pt]
\textit{(All three were doing the same work --- but only one truly understood the meaning of what they were doing.)}

Người thầy giáo đặt viên gạch cuối mỗi ngày --- không phải gạch đất nung --- mà là kiến thức, phẩm giá, và ước mơ trong trái tim học sinh.\\[4pt]
\textit{(The teacher lays the last brick each day --- not of fired clay --- but of knowledge, dignity, and dreams in students' hearts.)}

\end{truyen}
\begin{baihoc}
  Ý nghĩa công việc không nằm ở bản thân công việc --- mà ở cách bạn hiểu mình đang đóng góp vào điều gì lớn hơn.\\[4pt]
  \textit{(The meaning of work lies not in the work itself --- but in how you understand your contribution to something greater.)}
\end{baihoc}

\section{Ông Nội Dạy Cháu Câu Chuyện Hai Con Sói}
\begin{truyen}{Ông Nội Dạy Cháu Câu Chuyện Hai Con Sói}{Triết lý người Cherokee --- Bắc Mỹ}
\chuhoa{M}{ột} người Cherokee già nói với cháu: ``Trong lòng ta có một cuộc chiến --- cuộc chiến giữa hai con sói.''\\[4pt]
\textit{(An elderly Cherokee man told his grandchild: `Inside us there is a battle --- a battle between two wolves.')}

``Con sói thứ nhất là tức giận, ghen tị, đau khổ, tham lam, kiêu ngạo, tự ti, ích kỷ và dối trá.''\\[4pt]
\textit{(`The first wolf is anger, envy, sorrow, greed, arrogance, self-pity, guilt, and lies.')}

``Con sói thứ hai là vui vẻ, bình yên, tình yêu, hy vọng, khiêm tốn, lòng tốt, đồng cảm và sự thật.''\\[4pt]
\textit{(`The second wolf is joy, peace, love, hope, humility, kindness, empathy, and truth.')}

Đứa cháu hỏi: ``Con sói nào thắng?'' Ông nội trả lời: ``Con sói mà con cho ăn.''\\[4pt]
\textit{(The grandchild asked: `Which wolf wins?' The grandfather replied: `The one you feed.')}

\end{truyen}
\begin{baihoc}
  Mỗi ngày chúng ta nuôi dưỡng một trong hai con sói bằng những suy nghĩ và lựa chọn của mình --- và con sói nào mạnh hơn sẽ quyết định ta trở thành ai.\\[4pt]
  \textit{(Each day we feed one of the two wolves with our thoughts and choices --- and whichever wolf is stronger determines who we become.)}
\end{baihoc}

\section{Võ Thị Sáu --- Tuổi Mười Bảy Không Cúi Đầu}
\begin{truyen}{Võ Thị Sáu --- Tuổi Mười Bảy Không Cúi Đầu}{Võ Thị Sáu --- Việt Nam, 1933--1952}
\chuhoa{V}{õ} Thị Sáu tham gia kháng chiến từ năm 14 tuổi, bị Pháp bắt năm 16 tuổi, và bị xử tử tại Côn Đảo năm 19 tuổi.\\[4pt]
\textit{(Vo Thi Sau joined the resistance at 14, was captured by the French at 16, and was executed at Con Dao at 19.)}

Trước khi bị bắn, bọn lính bịt mắt cô; cô gạt tay lính ra: ``Tôi không sợ chết và không cần bịt mắt.''\\[4pt]
\textit{(Before being shot, soldiers tried to blindfold her; she pushed their hands away: `I am not afraid of death and do not need a blindfold.')}

Cô hát một bài hát yêu nước cho đến viên đạn cuối cùng.\\[4pt]
\textit{(She sang a patriotic song until the final bullet.)}

Cô ấy mang lại cho các thế hệ sau một bài học không thể học từ sách giáo khoa: lòng dũng cảm không phải là không sợ hãi --- mà là hành động dù đang sợ.\\[4pt]
\textit{(She gave subsequent generations a lesson no textbook can teach: courage is not the absence of fear --- it is acting despite being afraid.)}

\end{truyen}
\begin{baihoc}
  Tuổi trẻ không phải là lý do để chờ đợi --- lịch sử được viết bởi những người dám hành động, dù tuổi tác nào.\\[4pt]
  \textit{(Youth is not a reason to wait --- history is written by those who dare to act, at any age.)}
\end{baihoc}

\section{Thầy Nguyễn Ngọc Ký Viết Bằng Chân}
\begin{truyen}{Thầy Nguyễn Ngọc Ký Viết Bằng Chân}{Nguyễn Ngọc Ký --- Việt Nam, 1947--2022}
\chuhoa{N}{guyễn} Ngọc Ký bị liệt cả hai tay từ năm 4 tuổi; ông không thể cầm bút, không thể viết như bao bạn bè khác.\\[4pt]
\textit{(Nguyen Ngoc Ky was paralyzed in both arms from age 4; he could not hold a pen, could not write like other children.)}

Ông kiên trì học viết bằng chân --- ngón chân kẹp bút, miệt mài tập luyện cho đến khi chữ viết đẹp hơn nhiều bạn bình thường.\\[4pt]
\textit{(He persistently learned to write with his feet --- toes gripping the pen, training tirelessly until his handwriting was more beautiful than many able-bodied peers.)}

Bác Hồ đến thăm và tặng ông một chiếc bút --- phần thưởng ý nghĩa nhất một đứa trẻ từng nhận.\\[4pt]
\textit{(Uncle Ho visited and gave him a pen --- the most meaningful award a child had ever received.)}

Nguyễn Ngọc Ký sau này trở thành thầy giáo dạy văn, nhà văn với nhiều tác phẩm, và là nguồn cảm hứng cho hàng triệu học sinh Việt Nam.\\[4pt]
\textit{(Nguyen Ngoc Ky later became a literature teacher, author of many works, and an inspiration to millions of Vietnamese students.)}

\end{truyen}
\begin{baihoc}
  Không phải hoàn cảnh mà là ý chí mới quyết định điều ta có thể làm được --- thiếu đôi bàn tay chỉ là cái cớ, không phải rào cản.\\[4pt]
  \textit{(It is not circumstances but will that determines what we can do --- the absence of hands is an excuse, not a barrier.)}
\end{baihoc}
